\scanpage{451}
\mspage{228}

et on aura
\[
(23)\qquad X_g=X_{g,i}\amalg_{\partial\widetilde M}\widetilde M
\]
par un \unsure{homéom.} bien déterminé de $\partial\widetilde M$ sur une partie ouverte et fermée de $\partial X_{g,i}$. D'ailleurs, on aura une appl. continue can.
\[
(24)\qquad \unsure{U_0|\mathcal U_0}\longrightarrow X_{g,i}
\]
qui définit un homéom. de $X_{g,i}$ sur le \deleted{\unsure{variété}} \illegible{1} obtenu en \unsure{contractant} les $\underset{\inserted{\substack{\unsure{\parallel}\\(\partial D_j)\times\varepsilon_j}}}{\uline{\partial\Delta_j}}\subset U_0$ $(g-i+1\leq j\leq g)$ à l'aide des \unsure{projections} $\partial\Delta_j\simeq(\partial D_j)\times\varepsilon_j\longrightarrow D_j$.

D'autre part, chaque $\widetilde M_j$ $(1\leq j\leq g-i)$, rev. \inserted{des} orientations du ruban de Möbius, est isomorphe au cylindre $S^1\times I$, et son anti-involution canonique \deleted{s'identifie} est sans point fixe, et s'identifie à :
\[
(z,t)\longrightarrow(\unsure{-z}\deleted{\illegible{1}},-t)
\]
(où $S^1$ est identifié aux nbs complexes de module $1$, et $I:[-1;+1]$).

D'ailleurs, $X_0$ \inserted{\unsure{orienté|orientée},} avec \unsure{son} anti-involution de \unsure{revêtement} de $\unsure{Y_0|X_0}$ s'identifie à la sphère \deleted{\unsure{des}} ordinaire\inserted{, dans un espace vectoriel euclidien de dim $3$}, avec l'anti-involution $x\longmapsto-x$.

\scanpage{452}

En résumé :

\uline{Proposition} On peut obtenir (pour $0\leq i\leq g$ fixés) les couples $(X_g,\sigma_i)$, à \unsure{homéomorphismes} près, en prenant la sphère \deleted{\unsure{orientée}} \inserted{euclidienne orientée} $X_0$, avec son antipodisme standard $\tau$, en prenant un ens. de $2g$ disques \inserted{$\unsure{D'_j}\ (j\in J)$} \unsure{mutuellement} disjoints, stable par $\tau$, et une partie $I'\subset J/\tau=I$ ($I$ de card $g$) avec $\operatorname{card}I'=i$, et en \unsure{procédant} ainsi : pour \unsure{tt|$\forall$} $j\in I$, soit $\Delta_j=D'_j\cup D''_j$ la \inserted{réunion des deux} disques correspondants, et $S_j=\partial\Delta_j/\tau$, de sorte que $S_j$ est un cercle, soit \figure{[Encadrement ouvert en bas : $T_j=S_j\times\unsure{\mathbf I^{\varepsilon_j}|\mathbb I^{\varepsilon_j}}$.]} (où $\unsure{\mathbf I|\mathbb I}=[-1,+1]$), $\varepsilon_j=\{j'\in J\mid j'\text{ sur }\unsure{j|i}\}$\inserted{$\simeq$ ens. des \uline{orientations} de $S_j$}, $\unsure{\mathbf I^{\varepsilon_j}|\mathbb I^{\varepsilon_j}}$ tordu par $\varepsilon_j$, de sorte qu'on a un homéomorphisme can.

\margin{l'introduire plus tard}

\[
(25)\qquad \partial T_j\simeq\partial\Delta_j.
\]
\inserted{\unsure{i.e.} $\partial T\simeq\partial V_0$ \deleted{\unsure{id}} \illegible{1}}

\margin{de sorte que $X_g=\unsure{V_0|U_0}\amalg_{\partial\Delta}T$ (où $\Delta=\coprod_1^g\Delta_j$, $T=\coprod_1^gT_j$) est \uline{\unsure{orienté|orientée}}}

\margin{NB. $T_j$ est canoniquement orienté et l'isom (25) \uuline{respecte} \unsure{comme} \unsure{doit} \inserted{\uline{renverse}} l'\unsure{orientation}}

Ceci posé, l'antipodisme $\tau$ \unsure{induit} sur chaque $\unsure{\partial T_i|\partial T_j}$ \inserted{$\simeq S_j\times\varepsilon_j$}, \deleted{\illegible{1}}\inserted{l'}anti-involution \inserted{canonique} \deleted{\unsure{qui}} \inserted{\unsure{provenant de cette}} \deleted{\unsure{permutation}} de $\unsure{\partial T_i|\partial T_j}$, qui échange les deux composantes connexes.

Pour \inserted{\unsure{tt|$\forall$}} $j\in\unsure{I|I'}$, on prolonge $\dot\tau|\partial T_j$ à $T_j$ en deux anti-involutions $\dot\tau_j$, \inserted{$\dot\tau'_j$} de $T_j$, \deleted{et on} \deleted{\unsure{vérifie que}} de telle façon que

\scanpage{453}
\mspage{229}

a) \deleted{pour $j\in I'$}, $\dot\tau_j$ soit sans points fixes

b) \deleted{pour $j\in I'$}, $\dot\tau'_j$ ait un ens. de pts fixes homéomorphe à $S_j$ par la projection $T_j^{\unsure{\dot\tau_j|\dot\tau'_j}}\longrightarrow S_j$ (cf. plus bas \unsure{\}} pour des choix \unsure{particuliers} \unsure{et} explicites \unsure{\}}. --- on \unsure{voit} qu'on peut $\unsure{\overline{\mathrm m}|=}$ supposer $T_j^{\unsure{\dot\tau_j|\dot\tau'_j}}=S_j\times\{0\}$).

Soit, pour \unsure{tte|$\forall$} partie $I'\subset I$, $\tau_{I'}$ l'anti-involution de $X_g=\unsure{V_0|X_0}\amalg_{\partial T}T$ (où $V_0=X_0\setminus\bigcup_{j\in J}D'_j$) qui coïncide avec $\tau$ sur $V_0$, avec $\unsure{\tau'_j|\dot\tau'_j}$ sur $T_j$ pour $j\in I'$, avec $\unsure{\tau_j^0|\dot\tau_j}$ pour $j\in I\setminus I'$. Alors l'on a
\[
(26)\qquad X_g^{(\tau_{I'})}=\bigcup_{j\in I'}(S_j\times\{0\})
\]
donc si $\operatorname{Card}I'=i$ $(0\leq i\leq g)$, alors $\tau_{I'}$ est une $\sigma_i$.

Choix des $\dot\tau_j$, $\dot\tau'_j$. On choisit un isomorphisme $S_j\simeq\unsure{\mathbb U|\mathbf U}\overset{\unsure{\mathrm{df}}}{=}\{z\in\mathbb C\mid |z|=1\}$, d'où une orientation de $S_j$, et une bijection $\varepsilon_j\simeq\{\pm1\}$, \inserted{d'où $\unsure{\mathbf I^{\varepsilon_j}|\mathbb I^{\varepsilon_j}}\simeq\unsure{I|\mathbf I}=[-1,+1]$} et on prend ;
\[
(27)\qquad
\left\{\begin{aligned}
\dot\tau_{\unsure{j|g}}(z,t)&=\deleted{(\unsure{-z|z},-t)}\;(-z\exp\unsure{\pi it},-t)\\
\dot\tau'_{\unsure{j|g}}(z,t)&=(z\deleted{\unsure{\exp\pi it}},-t)
\end{aligned}\right.
\]

\margin{\uline{NB} Pour la définition de $\dot\tau'_j$, on n'a pas besoin du choix d'un iso $S_j\simeq\unsure{\mathbb U|\mathbf U}$}

\deleted{On a bien $\dot\tau^2$}
\[
\dot\tau|\partial T=\dot\tau'|\partial T=\tau|\partial T,
\]
On a bien $\dot\tau^2=\mathrm{id}$, $\dot\tau'{}^2=\mathrm{id}$

D'ailleurs

\scanpage{454}

On \unsure{note|notera} que
\[
(\dot\tau'\dot\tau)^2
=\dot\tau'\dot\tau\dot\tau'\dot\tau
=\unsure{\bigl((z,t)\longmapsto(z\exp 2i\pi t,t)\bigr)}
\]
ce n'est pas l'application identique --- \unsure{l'ens.} de ses pts fixes est \unsure{égal} à : \deleted{\unsure{$S_j$}}
\unsure{l'ens.} des $(z,t)$ tels que $t\in\{-1,0,+1\}$ --- i.e.
\[
T_j^{(\dot\tau'_j\dot\tau_j)^2}
=S_j\times[\partial\unsure{\mathbf I|\mathbb I}^{\varepsilon_j}\cup\{0\}]
\tag{28}
\]

\deleted{\unsure{On a}}

\deleted{Pour \unsure{pouvoir} écrire}

Soit \unsure{$\rho_j=$.}
\[
\dot\rho_j=\dot\tau'_j\dot\tau_j^{\deleted{\illegible{1}}}
\qquad [(z,t)\longmapsto(-z\exp\pi it,t)]
\tag{29}
\]
--- c'est un \unsure{autom.} de $T_j$ qui est l'\uline{identité} sur $\partial T_j$.
Et pour toute partie $I'$ de $I=J/\tau$, soit $\rho_{I'}$ l'\unsure{autom.} de $X_g$ qui est l'identité sur
$V_0=X_g-\bigcup_{j\in I-I'}T_j$,
et qui est $\rho_j$ sur $T_j$ pour $j\in I'$.

\margin{Posant $\rho_j=\rho_{\{j\}}$, (30) on aura simplement
$\rho_J=\prod_{j\in J}\rho_j$.
Sauf erreur, $\rho_j$ engendre le groupe
$\unsure{S\Gamma^{+}_{*}(T_j)}\simeq\mathbb Z$,
\unsure{donc les $\rho_j$ eng. un groupe} $\simeq\mathbb Z^I$.
\inserted{\unsure{(orienté sur}\illegible{1}\unsure{)}}}

\margin{\uline{NB} On a
$\rho_j=\tau^{\deleted{\illegible{1}}}_{\{j\}}\tau_{\varnothing}$.}

On aura donc
\[
\left\{
\begin{array}{l}
\tau_{I'}\tau_{I''}
=\rho_{I'_0}^{\deleted{\illegible{1}}}\rho_{I''_0}^{-1\,\deleted{\illegible{1}}}
\\[3pt]
\deleted{\unsure{\tau'_{I'}\tau'_{I''}}=\unsure{\rho_{I'_0}\rho_{I''_0}^{-1}}}
\\[3pt]
\deleted{\illegible{5}}
\end{array}
\right.
\tag{32}
\]
\deleted{\unsure{et le}}
\inserted{où
\deleted{$\unsure{J}=I'\cup I''-\unsure{I'\cap I''}$}
\deleted{$J=I'\cup I''-(I'\cap I'')$}
$I'_0=I'\setminus I'\cap I''$, $I''_0=I''\setminus I'\cap I''$.}

d'autre part on aura évidemment
\[
[\rho_J,\rho_K]=1\qquad\text{pour }J,K\subset I.
\tag{33}
\]

\uline{Remarque.} Au lieu d'un \unsure{isom.} $S_j\simeq\mathbb U$,
supposons donné plutôt sur

\scanpage{455}
\mspage{230}

$S_j$ une structure de torseur sous $\mathbb U^{\varepsilon_j}$
($\mathbb U$ tordu par $\varepsilon_j$, \unsure{pour|pris} : l'\unsure{autom.} d'ordre 2
$z\longmapsto z^{-1}=\bar z$ de $\mathbb U$) On peut donc définir
\[
\exp:\underset{\underset{\unsure{\mathbf I|\mathbb I}^{\varepsilon_j}}{\cup}}{\mathbb R^{\varepsilon_j}}
\longrightarrow\mathbb U^{\varepsilon_j}
\tag{36}
\]
de façon évidente, d'où des \deleted{applications}
\inserted{\uline{anti-involutions}}
$\dot\tau_j,\dot\tau'_j:T_j\xrightarrow{\sim}T_j$ par les formules (27).

Si p.ex. on choisit des disques $D'_j$ tels que leurs bords soient des \uline{cercles}
\unsure{euclidiens}, alors il y a sur chaque \deleted{\unsure{$S_j$}} $\partial D'_j$
une structure de torseur sous un groupe $\mathbb U^{\varepsilon_j}$,
\unsure{invariant|invariante} par antipodisme, et qui passe au quotient.
Dès lors \unsure{tout} \unsure{autom.} de $(X_0,(D_j))$ qui respecte cette
\deleted{\unsure{montre que le}}
structure supplémentaire de torseurs --- et évidemment tt \unsure{autom.} qui respecte
la structure \uline{métrique}, \deleted{\unsure{respecte}} opère sur $X_g$ en commutant
au système des $\tau_J,\tau'_J$ en un sens évident.

\unsure{Égal\textsuperscript{t}}, si on \deleted{prend} \inserted{utilise} sur $X$
des structures \uline{conformes} seules, et si on

\scanpage{456}

choisit dans chaque $D'_j$ un ``centre'' $a'_j$, de façon compatible avec l'involution,
alors le choix des $a'_j\in\mathrm{Int}(D'_j)$ définit une structure de torseur sur
$\partial D'_j$\unsure{\textsuperscript{*}}, et ces structures sont invariantes par les
transformations conformes qui respectent \unsure{l'ens.} des pts $a'_j$.

\margin{\uline{NB} On ne suppose plus ici que $D_j$ soit un cercle,
\unsure{seulement} que $\partial D_j$ pas trop \deleted{\unsure{irrégulier}}.
Si $\partial D_j$ est un cercle, cette définition coïncide avec la précédente
\unsure{ssi} \unsure{$a_j$|$a'_j$} est le centre \unsure{du|de ce} cercle\ldots}

\uline{Re-NB.} Pour définir $\tau_j$ sur $T_j$, il suffit de \unsure{nettement}
moins de données que d'une structure de torseur topologique sur $T_j$.
Écrivons \inserted{pour $(z,t)\in S_j\times\unsure{\mathbf I|\mathbb I}^{\varepsilon_j}$}
\[
\tau_j(z,t)=(u_t(z),-t)
\tag{37}
\]
où $u_t:S_j\xrightarrow{\sim}S_j$ est un \unsure{homéom.} dépendant continûment de $t$,
écrivant que $\tau_j|\partial T_j=\tau|\partial T_j$, on trouve la condition
\[
\text{a)}\quad u_t=\deleted{\unsure{u_1=}}\mathrm{id}
\quad\text{si }t\in\partial\unsure{\mathbf I|\mathbb I}^{\varepsilon_j}\simeq\varepsilon_j
\]

\margin{\unsure{ça} s'écrit, si $S_j$ est orienté, $u_1=u_{-1}=\unsure{0}$.}

écrivant que $\tau_j^2=\mathrm{id}$ on trouve la condition
\[
\text{b)}\quad u_{-t}=u_t^{-1}
\qquad(t\in\unsure{\mathbf I|\mathbb I}^{\varepsilon_j})
\]
et écrivant que $T_j^{\tau_j}=\varnothing$ on trouve
\[
\text{c)}\quad u_0\text{ sans points fixes}.
\]

Si une orientation est choisie, les $t\longmapsto u_t$ satisfaisant a) b) c)
correspondent aux applications
$[0,1]\xrightarrow{t\mapsto u_t}\deleted{\unsure{\mathrm{Homéom}}}\inserted{\unsure{\mathrm{Aut}}}(\unsure{S_j|S_1})$
telles que $u_1=\mathrm{id}$, $u_0$ ss pt fixe
\deleted{\illegible{2}}\inserted{\uuline{et d'ordre 2}}.
(Le cas \unsure{envisagé} plus haut est celui où
$\unsure{*}t\longmapsto u_{1-t}$ provient d'une représentation continue
$\mathbb R\longrightarrow\mathrm{Aut}(\unsure{S_j|S_1})$\ldots)

\scanpage{457}
\mspage{231}

\uline{Remarques} On peut se proposer de déterminer la structure de toutes les anti-involutions \inserted{$\tau$} sur un cylindre \inserted{orienté} $T\simeq S\times\unsure{\mathbf I^{\varepsilon}|\mathbb I^{\varepsilon}}$ $(\varepsilon=\unsure{\operatorname{or}(S)|\operatorname{Or}(S)})$ qui \deleted{\illegible{1}} n'ont pas de pt fixe sur le bord --- ce qui implique déjà qu'\deleted{elles}\inserted{$\tau$} \unsure{permute} les deux composantes connexes du bord. Plus généralement, les anti-involutions \inserted{$\tau$} d'une \deleted{\unsure{variété}}\inserted{surface} \uline{à bord} orientée $\unsure{X}$, n'ayant pas de pt fixe sur $\partial\unsure{X}$, correspondent aux variétés à bord $\unsure{Y}$ \unsure{munies} d'une partie à la fois \unsure{ouverte} et \unsure{fermée} $(\partial Y)'$ de $\partial Y$ --- on associe à une telle $\unsure{(Y,(\partial Y)')}$ son \uline{``\unsure{doublement}'' orienté} relativement à $(\partial Y)'$ --- à $X,\sigma$ correspond $(X/\sigma,\operatorname{Im}X^\sigma\longrightarrow X/\sigma)$. Si $X$ est connexe, \inserted{compact,} $Y$ est \inserted{compacte,} \uline{non orientable}, supposons que \unsure{son type} $(\gamma,j)$, \deleted{soit} et \deleted{\unsure{card} $\pi_0$} \unsure{soit} \inserted{$0\leq i\leq j$, et,} $\unsure{i|i'}=\operatorname{card}(\pi_0((\partial Y)'))$\inserted{$={}\operatorname{card}\pi_0(X^\sigma)$}. Donc $\chi(Y)=1-(\gamma+j)$, et on sait de \unsure{suite|sorte} que \deleted{$\chi(X)=\chi_!(\operatorname{Int}(X))$} (car le $\chi$ d'un cercle est nul)
\[
\begin{split}
\chi(X)&=\chi_!\left(\underbrace{X\setminus(\unsure{X\mid\partial Y})}_{\unsure{X\mid\operatorname{Int}(Y)}}\right)\\
&=2\chi_!(\operatorname{Int}Y)=2\chi(Y)=2(1-(\gamma+j))
\end{split}
\]
donc si $X$ est de \deleted{genre}\inserted{type} $(g,\nu)$, on aura \inserted{$\chi(X)=2-2g\unsure{-|+}\nu$}
\[
\deleted{\unsure{g+\nu}}\quad 2-2g\unsure{-|+}\nu=2(1-(\gamma+j))
\]
i.e.
\[
(38)\qquad g+\frac{\nu}{2}=\gamma+j
\]

\margin{pas vrai ! Si $Y$ orientable de type $(\gamma,j)$ (\unsure{même|avec} $0\leq i\leq j$) on aura
\[
(39)\qquad\left\{\begin{array}{l}
g+\dfrac{\nu}{2}=2\gamma+j-1\\
\nu=2(j-i)
\end{array}\right.
\]
i.e. $i=j-\dfrac{\nu}{2}$ (comme dans le cas ci-contre)}

\scanpage{458}

(cela exige que $\nu$ soit pair, ce qui était évident à priori, car $\sigma$ doit permuter les éléments de $\pi_0(\partial X)$ entre eux sans \inserted{y} avoir de pt fixe\ldots)
\deleted{Dans le cas actuel $g=1,\nu=2$ on \unsure{trouve|trouvera} donc}
\[
\gamma+i=2
\]
\deleted{d'où les possibilités}
\deleted{$\gamma=2,\ i=0$}
\inserted{Mais j'ai oublié de noter que $\nu$ est déterminé en fonction de $(\gamma,j;i\inserted{={}\unsure{\operatorname{card}((\partial Y)')}})$ par}
\deleted{$\nu=2(j-i)+\unsure{i^{\varepsilon}|i}=2j-i$}
\deleted{$\unsure{\gamma|\nu}=2(j-i)$ \unsure{donc}}
\deleted{\unsure{doit avoir}}
\[
(39)\qquad \nu=2(j-i)
\]
i.e. \uline{$i=j-\dfrac{\nu}{2}$}.

Donc il faut ici (pour $g=0,\nu=2$) chercher $(\gamma,j;i)$, avec $\gamma\in\unsure{\mathbb N|N},\ 0\unsure{\leq|<}i\leq j\in\unsure{\mathbb N|N}$, tels que l'on ait $\nu=2(j-i)$ i.e. $j-i=1$ i.e. $i=j-1$, et $\gamma+j=1$, ce qui donne \unsure{la|les} \inserted{\unsure{seule}} \unsure{possibilité|possibilités} (comme $i$ doit être $\geq0$ donc $j\geq1$)

a) $\gamma=0,\ j=1,\ i=0$ \unsure{:|.}

\deleted{b) $\unsure{\gamma=1|\gamma=0},\ \unsure{j=0|j=1},\ i=0$.}

Le cas a) correspond au cas où \deleted{\illegible{1}} $T^\sigma\unsure{=|\neq}\varnothing$, \deleted{il \illegible{2} \unsure{nécessairement} une composante connexe, et} se déduit du plan projectif \unsure{réel|vide} en y faisant \deleted{deux}\inserted{un} \unsure{trou|trous}, \inserted{\unsure{à bord (d'où ruban Möbius)}} et en \deleted{bordant l'un d'eux par}\inserted{prenant le \deleted{\unsure{rev. orienté}} doublement orienté}. \deleted{relativement à l'un de ces trous.}

Le cas où $T^\sigma\neq\varnothing$ \unsure{donnera} nécessairement un quotient $Y$ \uline{\unsure{orienté|orientable}}, on doit avoir pour son type $\unsure{(\gamma,j)}$ (\unsure{car|comme} $j\geq1$), le seul cas :

\scanpage{459}
\mspage{232}

b) $\gamma=0,\ j=2,\ i=1$ (quotient \uline{orienté})

déduit de la sphère à deux trous i.e. du cylindre, en prenant le doublement orienté \unsure{p.r.} à \uline{un} des trous. \deleted{C'est bien les}

C'est bien les deux cas obtenus \unsure{resp.\textsuperscript{t}} par $\dot\tau$ et $\dot\tau'$ dans les formules (27).

Je voudrais maintenant construire une situation (sous les conditions de la proposition) mettant en évidence un \unsure{maximum|mécanisme} de \unsure{symétries} --- il est vrai que l'on pourrait \unsure{travailler} avec \deleted{\unsure{le} \illegible{1} \unsure{des}} les \unsure{autom.} de $X_0$ \unsure{commutant} à l'antipodisme\inserted{$\tau$}, \inserted{\unsure{laissant|conservant} \unsure{l'ens} des $D'_j$,} et respectant (disons) des structures de \unsure{torseurs|torseur} topologiques sur les $\partial D'_j$ --- ou ce qui revient \illegible{1} au $\overline{\mathrm m}$ (à \unsure{indétermination} de multiplication par $\unsure{\tau|2}$ près) les \unsure{autom.} de $Y_0$ qui \deleted{\illegible{1}} \unsure{invarient|invariant} $D=\bigcup_{i\in I}D_i$, et \unsure{respectent|respectant} des structures de \unsure{torseur|torseurs} sur \deleted{les bords} les comp. conn. $\partial D_i$ de $\partial D$. On prévoit que (\unsure{travaillant|Au moins} mod isotopie) on \unsure{obtient|aura} un groupe

\margin{non, \uline{sans} \unsure{indétermination}, \unsure{un relèvement|en relevant} \unsure{défini} de façon à \unsure{respecter} les \unsure{orientations|orientation}\ldots}

\scanpage{460}

qui sera voisin d'un groupe de tresses, et sans doute calculable sans grand mal --- et en \inserted{\unsure{le}} mettant ens. avec le groupe engendré par les opérations précédentes, on trouvera peut être un \unsure{dévissage} pour \unsure{engendrer} p.ex. \unsure{$\Pi_g$|$\Gamma_g$} par générateurs \inserted{(anti-involutifs ?)} et relations.

Considérons le \deleted{\illegible{1}} groupe $G$ de $A_g=\operatorname{Aut}X_g$ engendré par les \deleted{\illegible{1}}\inserted{$\tau_{I'}$} $(I'\unsure{\in|\subset} I)$, \deleted{les} Soit $\mathcal H$ \deleted{commutateurs de générateurs sont dans} le sous-groupe engendré par les $\rho_j$ $(j\in I)$. On a
\[
\begin{aligned}
\tau_{I'}\rho_j\tau_{I'}
&=\deleted{\unsure{(\tau_{I'}(\tau_{\{j\}}\tau_{\varnothing})\tau_{I'})}}
\tau_{I'}(\tau_{\{j\}}\tau_{\varnothing})\tau_{I'}\\
&=(\tau_{I'}\tau_{\{j\}})
\underbrace{(\tau_{\varnothing}\tau_{I'})}_{\deleted{\illegible{1}}\unsure{\rho_{I'}}}
\in\mathcal H
\end{aligned}
\]
par les formules (31), donc $\mathcal H$ est un sous-groupe invariant.

Les formules (32) montrent que $G/\mathcal H$ est un groupe \uline{commutatif}, \deleted{\unsure{isomorphe à}}\inserted{et $\overline{\mathrm m}$ qu'il est \unsure{isomorphe à}} $\pm1$ par le caractère d'orientation (tous les $\tau_{I'}$ sont égaux mod $\mathcal H$).
